\chapter{Perfect and finite duality}
\label{ch:derived-duality}

All rings in the affine part of this chapter are commutative.
We distinguish internal Hom over a ring from the Hom groups in its
derived category. For a map \(A\to B\) we keep the \(B\)-action on
the right adjoint explicitly.

\section{Perfect complexes}

\begin{theorem}[Evaluation and duals of perfect complexes]\label{res:U10}
Let \(P\in D(R)\) be represented by a bounded complex of finitely
generated projective modules. Define
\(P^\vee=\mathbf R\mathrm{Hom}_R(P,R)\).
Then \(P^\vee\) is perfect, the natural evaluation comparison
\[
 Y\otimes_R^{\mathbf L}P^\vee
     \longrightarrow\mathbf R\mathrm{Hom}_R(P,Y)
\]
is an isomorphism for every \(Y\), and \(P\) is dualizable
with dual \(P^\vee\). Biduality \(P\to(P^\vee)^\vee\) is an
isomorphism.
\end{theorem}
\begin{proof}
Use the bounded finite projective representative for \(P\).
It is K-projective by \cite[R03]{ha-core} and K-flat by
Proposition~\ref{res:U03}. The ordinary dual
\(D=\underline{\mathrm{Hom}}_R(P,R)\) is bounded finite
projective. The comparison in total degree \(q\) sends
\(y\otimes f\) to the map \(x\mapsto y f(x)\) for
\(|y|+|f|=q\). Expanding the differential gives
\(d y\, f-(-1)^{|y|+|f|} y f d_P\), as required.
It is a termwise isomorphism, since each degree of \(P\)
is finitely generated projective and the diagonals involving
\(P\) are finite. The standard evaluation \(D\otimes_R P\to R\)
sends \(f\otimes x\) to \(f(x)\). Applying the isomorphism just
proved with \(Y=P\), the identity of \(P\) corresponds to a
degree-zero cycle in \(P\otimes D\). This defines coevaluation
\(R\to P\otimes D\). In each degree choose a finite dual basis.
The cycle is the sum of the dual basis tensors over the
finitely many degrees. Both triangle identities now follow
by the dual basis expansions in each degree, with evaluation
on the adjacent \(D,P\) factors. In particular the identities
hold already on the complexes. The bidual map is the
degreewise finite projective double dual, with the symmetry
sign \(x\mapsto[f\mapsto(-1)^{|x||f|}f(x)]\).
It is a chain map by the Hom differential and a
termwise isomorphism.
\end{proof}

\begin{proposition}[Projection through a perfect factor and base change]\label{res:U11}
If \(P\) is perfect, the natural comparison
\[
 P\otimes_R^{\mathbf L}\mathbf R\mathrm{Hom}_R(X,Y)
 \longrightarrow
 \mathbf R\mathrm{Hom}_R(X,P\otimes_R^{\mathbf L}Y)
\]
is an isomorphism. Also, for any \(R\to S\) the natural
change-of-rings map
\[
 S\otimes_R^{\mathbf L}\mathbf R\mathrm{Hom}_R(P,Y)
  \longrightarrow
 \mathbf R\mathrm{Hom}_S(S\otimes_R^{\mathbf L}P,
                          S\otimes_R^{\mathbf L}Y)
\]
is an isomorphism.
\end{proposition}
\begin{proof}
For the first assertion test by morphisms from any \(Z\).
Using Theorem~\ref{res:U10}, morphisms from \(Z\) to the
left side are morphisms from \(P^\vee\otimes^{\mathbf L}Z\)
to \(\mathbf R\mathrm{Hom}_R(X,Y)\). The closed adjunction
in Theorem~\ref{res:U08} and associativity identify them with
morphisms from \(P^\vee\otimes^{\mathbf L}(Z\otimes^{\mathbf L}X)\)
to \(Y\). Applying the same duality a second time and
then the closed adjunction gives morphisms from \(Z\) to the
right side. The composite is the canonical tensor-evaluation
comparison, since each step takes adjoints under evaluation
or coevaluation. The isomorphism follows by Yoneda.
For base change, extend scalars on a bounded finite
projective representative for \(P\) and its ordinary dual.
Finite projective duals commute termwise with extension,
and the comparison is compatible with evaluation.
Proposition~\ref{res:U09} transports the evaluation comparison
of Theorem~\ref{res:U10} to the evaluation comparison over \(S\).
This proves the assertion even if \(S\) is not flat over \(R\).
\end{proof}

\begin{example}[Termwise free need not be K-flat]\label{res:U12}
Let \(k\) be a field and \(R=k[\epsilon]/(\epsilon^2)\).
The unbounded complex \(A\) having \(R\) in every degree and
differential multiplication by \(\epsilon\) is acyclic.
It is neither K-projective nor K-flat.
\end{example}
\begin{proof}
Its cycles and boundaries are both \(\epsilon R\) in every degree.
It cannot be K-projective, since its identity is not null-homotopic.
In any proposed homotopy \(d h+h d\) has image in \(\epsilon R\).
For the tensor complex \(A\otimes A\), use in degree \(n\) the
finitely supported basis \(e_i^{(n)}\), indexed by \((i,n-i)\).
Then
\[
 d(e_i^{(n)})=\epsilon e_{i+1}^{(n+1)}
                 +(-1)^i\epsilon e_i^{(n+1)}.
\]
The element \(\epsilon e_0^{(n)}\) is a cycle. If it were a boundary,
a nonzero finitely supported sequence \((a_i)_i\) over \(k\)
would have image \(e_0^{(n)}\) under
\(e_i^{(n-1)}\mapsto e_{i+1}^{(n)}+(-1)^i e_i^{(n)}\).
For such a sequence the least and greatest image indices are its
least index and one more than its greatest index. They are
distinct, a contradiction. This shows directly that \(A\otimes A\)
is not acyclic.
\end{proof}

\begin{example}[An infinite source in Hom base change]\label{res:U13}
Let \(R=k\) be a field, \(S=k[t]\), \(P=\bigoplus_{n\geq0}k\),
and \(Y=k\), all in degree zero. The canonical Hom base-change
map for \(P,Y\) is not surjective.
\end{example}
\begin{proof}
The map is \(k[t]\otimes_k\prod_{n\geq0}k\to\prod_{n\geq0}k[t]\),
by ordinary projectivity and flatness. Any element of its image
has a common bound on the degrees of its entries, since it is a
finite sum of polynomial multiples of scalar sequences.
The sequence \((t^n)_{n\geq0}\) violates that condition.
\end{proof}

\section{Finite affine adjunction and traces}

\begin{theorem}[Derived coinduction and its trace]\label{res:U14}
For a ring map \(A\to B\) let \(f_*:D(B)\to D(A)\) be
restriction of scalars. It has a right adjoint \(f^!\).
If \(Y\to I\) is a K-injective replacement over \(A\), it is computed
by \(f^!Y=\underline{\mathrm{Hom}}_A(B,I)\) with
\((b\cdot\varphi)(b')=\varphi(bb')\).
The counit \(f_*f^!Y\to Y\) is evaluation at \(1\).
Under successive ring maps these right adjoints compose and
the counits compose by successive evaluations at \(1\).
\end{theorem}
\begin{proof}
For any \(B\)-complex \(M\) there is an isomorphism of Hom complexes
\[
 \underline{\mathrm{Hom}}_B
  (M,\underline{\mathrm{Hom}}_A(B,I))
 \longrightarrow \underline{\mathrm{Hom}}_A(f_*M,I),
 \qquad u\longmapsto [m\mapsto u(m)(1)].
\]
Its inverse sends \(v\) to \(m\mapsto[b\mapsto v(bm)]\).
Since \(B\) is in degree zero both are chain maps by the ordinary
Hom differential. Exactness of restriction and K-injectivity
of \(I\) show that \(\underline{\mathrm{Hom}}_A(B,I)\)
is K-injective over \(B\). The displayed isomorphism of degree-zero
cohomologies therefore gives the derived adjunction by
\cite[R04]{ha-core}. It identifies the image of the identity
with evaluation at \(1\), hence gives the counit.
For \(A\to B\to C\) apply the same expression with \(M=C[0]\).
It gives a \(C\)-linear isomorphism
\[
 \underline{\mathrm{Hom}}_B(C,\underline{\mathrm{Hom}}_A(B,I))
 \cong\underline{\mathrm{Hom}}_A(C,I).
\]
The inner model on the left is K-injective, so the left side indeed
computes the composite. Since evaluation at \(1_C\) in the
image means evaluating first at \(1_C\) and then at \(1_B\),
the trace statement follows.
\end{proof}

\begin{proposition}[Finite relative dual under a perfectness hypothesis]\label{res:U15}
If \(B\) is perfect as an \(A\)-complex, put
\(D_{B/A}=f^! A\in D(B)\).
For every \(Y\in D(A)\) the natural \(B\)-linear map
\[
 D_{B/A}\otimes_B^{\mathbf L}(B\otimes_A^{\mathbf L}Y)
     \longrightarrow f^!Y
\]
is an isomorphism.
\end{proposition}
\begin{proof}
Choose a K-injective representative \(I\) of \(A\) and a
semifree representative \(P_Y\) for \(Y\).
The left side is represented by
\(\underline{\mathrm{Hom}}_A(B,I)\otimes_A P_Y\)
because \(B\otimes_A P_Y\) is K-flat over \(B\) by the proof
of Proposition~\ref{res:U09}. Its canonical map to
\(\underline{\mathrm{Hom}}_A(B,I\otimes_A P_Y)\)
sends a pure tensor \(\varphi\otimes p\) to
\([b\mapsto\varphi(b)\otimes p]\).
The corresponding map for a general complex \(Q\) in place of
\(B\) uses \((-1)^{|p||q|}\varphi(q)\otimes p\) on \(q\in Q\).
This is a chain map by the two differentials of
Definition~\ref{res:U01}. Compose with Hom of a
K-injective replacement of \(I\otimes_A P_Y\), which represents
\(Y\). To see this map is a quasi-isomorphism, forget to
\(A\) and resolve \(B\) by a bounded finite projective complex
\(Q\). Replacing \(B\) by \(Q\) in the two Hom complexes with
K-injective targets is a quasi-isomorphism by
Proposition~\ref{res:U03}. Tensoring the comparison for \(I\)
with \(P_Y\) preserves this property. For the \(Q\)-model,
ordinary tensor evaluation is an isomorphism by
Theorem~\ref{res:U10}, and Hom from \(Q\) to the chosen
replacement is a quasi-isomorphism by its K-projectivity.
Commutativity of the comparison squares proves the original map
is a quasi-isomorphism. This construction retains \(B\)-linearity
because the map on the unreplaced \(B\)-models does.
\end{proof}

\begin{proposition}[Flat base change for a finite algebra]\label{res:U16}
Let \(A\) be Noetherian, \(B\) be a finite \(A\)-algebra,
and \(A\to A'\) be flat. Set \(B'=A'\otimes_A B\).
For \(Y\in D^+(A)\), coinduction satisfies the \(B'\)-linear
base-change isomorphism
\[
 B'\otimes_B^{\mathbf L}f^!Y
   \ \cong\ f'^!(A'\otimes_A^{\mathbf L}Y).
\]
It identifies the two evaluation-at-\(1\) traces.
\end{proposition}
\begin{proof}
Take a bounded-below termwise injective representative \(I\) of
\(Y\) and a bounded-below injective replacement \(I'\) of
\(A'\otimes_A I\), using \cite[R01, R03]{ha-core}.
Since \(B\) is finitely presented as an \(A\)-module, flat
extension commutes with \(\mathrm{Hom}_A(B,I^q)\) in each
degree, by taking the kernel of Hom applied to a finite
presentation. Thus the comparison is the ordinary Hom
of \(B'\) with \(A'\otimes_A I\to I'\) on these models.
The \(B\)-base change is underived on every term because
\(B'\) is \(B\)-flat.

For the quasi-isomorphism assertion resolve \(B\) over \(A\)
by a resolution \(P_B\) in nonpositive degrees of finite free
modules. Such a resolution is available by iterating finite
free surjections onto successive finite kernels over the
Noetherian ring. Its scalar extension resolves \(B'\)
by flatness. In each total degree the products in
\(\underline{\mathrm{Hom}}_A(P_B,I)\) are finite because
\(P_B\) is bounded above and \(I\) bounded below. Hence
ordinary Hom on this model commutes with scalar extension
termwise. On the base-changed model, Hom from the
bounded-above projective resolution to \(A'\otimes I\to I'\)
is a quasi-isomorphism by \cite[R03]{ha-core}.
Comparing the two projective models with Hom against
the injective targets proves the claim, since flat tensor
preserves the comparison on the \(A\)-side.
For trace, evaluation at \(1_B\) commutes with change of
scalars and comparison of the target complexes. That is the
counit of Theorem~\ref{res:U14}.
\end{proof}

\begin{corollary}[A quotient by a regular sequence]\label{res:U17}
Let \(A\) be Noetherian, \(J=(a_1,\ldots,a_c)\) be generated
by an \(A\)-regular sequence, and \(B=A/J\).
Then
\[
  f^!Y\ \cong\
  \bigl(\textstyle\bigwedge\nolimits^c_B(J/J^2)\bigr)^\vee[-c]
       \otimes_B^{\mathbf L}(B\otimes_A^{\mathbf L}Y).
\]
Here the dual on the determinant line is its ordinary module dual.
\end{corollary}
\begin{proof}
The Koszul complex on the \(a_i\) resolves \(B\) in degrees
\(-c,\ldots,0\) by \cite[RK-07]{ca-core}. Its graded free
algebra is the exterior algebra on a free module \(F\) with
basis sent to the \(a_i\). In degree \(i\), identify the dual
Koszul term with \(\bigwedge^{c-i}F\otimes(\bigwedge^c F)^\vee\)
by \(z\otimes\lambda\mapsto[x\mapsto\lambda(z\wedge x)]\).
For \(x\in\bigwedge^{i+1}F\) the identity
\(d(z\wedge x)=d z\wedge x+(-1)^{c-i}z\wedge d x=0\)
shows that the Hom differential acts on \(z\) as
\((-1)^{i+1}(-1)^{c-i+1}d z=(-1)^c d z\).
This is the shift by \(-c\). Thus
its cohomology vanishes outside degree \(c\), where it is
\(B\otimes_A(\bigwedge^c F)^\vee\). The action of \(B\) on
this Ext group is the coinduction action. To check it, lift an
element of \(B\) to \(A\). Multiplication on the resolution
by the lift is a chain map lifting multiplication on \(B\),
and induces the ordinary scalar action on cohomology.
If the lift differs by an element of \(J\), the corresponding
multiplication is null-homotopic on the Koszul resolution by
wedge with the corresponding degree \(-1\) element.

The basis of \(F\otimes_A B\) identifies with the classes of
the \(a_i\) in \(J/J^2\). To justify independence, whenever
\(\sum_i b_i a_i\) lies in \(J^2\), subtract a relation with
coefficients in \(J\). The remaining exact relation is in
the image of the degree-two Koszul differential and hence its
coefficients lie in \(J\). Surjectivity follows from generation.
A change of regular generators induces the corresponding
matrix on this conormal module. Lifting its coefficients
gives a comparison of Koszul complexes whose exterior
powers identify the top Ext via the dual determinant of
that same matrix. It follows that the top Ext identification
is canonical. A complex of \(B\)-modules with this unique
nonzero cohomology is canonically represented in \(D(B)\)
by that module in degree \(c\), by good truncation.
The Koszul resolution makes \(B\) perfect over \(A\), so
Proposition~\ref{res:U15} gives the formula for \(Y\).
\end{proof}

\section{Comparison with geometric and local duality}

\begin{proposition}[The finite morphism between affines]\label{res:U18}
Let \(f:\mathrm{Spec}(B)\to\mathrm{Spec}(A)\) be finite.
On the derived categories of complexes of quasi-coherent
modules on these affine schemes, \(f_*\), \(\mathbf Lf^*\),
and the right adjoint to \(f_*\) correspond respectively to
restriction, derived scalar extension, and the coinduction
in Theorem~\ref{res:U14}. The finite trace in this adjunction
corresponds to evaluation at \(1\).
\end{proposition}
\begin{proof}
We recall the needed exact affine correspondence.
For an \(A\)-module \(M\), the module sheaf \(\widetilde M\)
has sections \(M_a\) on \(D(a)\).
Conversely for a quasi-coherent \(\mathcal F\) choose a finite
distinguished affine cover on which it is a module sheaf.
Writing the sheaf equalizer for sections over this cover
and then localizing at \(a\) gives the equalizer for the
induced cover of \(D(a)\), since localization is exact
and the products in this equalizer are finite. Hence
\(\Gamma(X,\mathcal F)_a=\Gamma(D(a),\mathcal F)\).
The canonical map \(\widetilde{\Gamma(X,\mathcal F)}\to
\mathcal F\) is an isomorphism on the distinguished basis.
This equivalence is exact because the module sheaf functor
is exact on each stalk by exactness of localization. It
induces an exact equivalence on derived categories by
termwise application and localization. Preimages of \(D(a)\)
under \(f\) are \(D(a)\) in \(\mathrm{Spec}(B)\), where \(a\)
is viewed in \(B\). Direct image thus corresponds to
restriction of scalars, while ordinary pullback on affines
is extension of scalars by its tensor description.
Apply Proposition~\ref{res:U09} and
Theorem~\ref{res:U14}. The trace is their adjunction counit,
hence the requested evaluation.
\end{proof}

\begin{theorem}[Comparison with complete local duality]\label{res:U19}
Let \((R,\mathfrak m,k)\) be a complete Noetherian local ring.
Choose a surjective local map \(f:S\to R\) from a complete
regular local ring \((S,\mathfrak n,k)\) of dimension \(d\),
with the same residue field. Then
\[
 D_R=f^!(S[d])=\mathbf R\mathrm{Hom}_S(R,S[d])
\]
is the normalized dualizing complex of
\cite[D05]{ca-topics}. For finite \(R\)-modules \(M\)
there are natural isomorphisms
\[
 \begin{split}
 \mathbf R\mathrm{Hom}_R(M,D_R)
      &\cong \mathbf R\mathrm{Hom}_S(M,S[d]),\\
 H^{-q}\mathbf R\mathrm{Hom}_R(M,D_R)
      &\cong \mathrm{Hom}_R(H^q_{\mathfrak m}(M),E_R(k)).
 \end{split}
\]
In the first line the \(R\)-action on the right is by
precomposition on \(M\). The finite trace is the same
evaluation as in Theorem~\ref{res:U14}.
\end{theorem}
\begin{proof}
The regular cover and the identification of the normalized
dualizing complex are \cite[C04, D05]{ca-topics}.
The first isomorphism is the complex-level adjunction of
Theorem~\ref{res:U14}. Multiplication on the \(R\)-factor
there gives multiplication by precomposition on \(M\),
as seen from the inverse formula for \(v\).

Here is a direct comparison with the local cohomology
proof. Put \(E_S=E_S(k)\) and let \(C_{\mathfrak n}\)
be the finite flat augmented \v{C}ech complex for
generators of \(\mathfrak n\). By
\cite[L05, D01, M03]{ca-topics}, \(C_{\mathfrak n}\)
represents \(E_S[-d]\) and
\(\mathrm{End}_S(E_S)=S\).
The complex \(W=\underline{\mathrm{Hom}}_S(C_{\mathfrak n},E_S)\)
is bounded and termwise injective, since Hom from a flat
module into an injective module is injective by tensor--Hom
adjunction and flatness. Exactness of dualization into
\(E_S\) shows that \(W\) has only one nonzero cohomology
module, the copy of \(S\) in degree \(-d\). A cycle representing
its generator \(1\) gives a quasi-isomorphism \(S[d]\to W\).
Thus it may be used for the K-injective model in
Theorem~\ref{res:U14}. Since the images of the generators
generate \(\mathfrak m\), scalar extension identifies
\(R\otimes_S C_{\mathfrak n}=C_{\mathfrak m}\)
by localization. Coinduction identifies
\(E_R(k)\cong\mathrm{Hom}_S(R,E_S)\) by
\cite[D04]{ca-topics}. The complex adjunction yields
\[
 \underline{\mathrm{Hom}}_S(R,W)
    \cong \underline{\mathrm{Hom}}_R(C_{\mathfrak m},E_R).
\]
This is a K-injective \(R\)-model by Theorem~\ref{res:U14}.
Applying \(\underline{\mathrm{Hom}}_R(M,-)\) gives
\(\underline{\mathrm{Hom}}_R(M\otimes_R C_{\mathfrak m},E_R)\).
Taking degree \(-q\) cohomology gives the displayed Matlis
dual, since Hom into an injective module is exact and the
tensor computes local cohomology. The comparison agrees
with \cite[D06]{ca-topics}, using the same bounded
injective model and the same change of scalars. The trace
on the coinduced model is evaluation at \(1_R\) as claimed.
\end{proof}

\begin{corollary}[Completion and the Cohen--Macaulay form]\label{res:U20}
For a Noetherian local ring \(R\) and a finite \(R\)-module \(M\)
there is an isomorphism
\[
 \mathrm{Hom}_R(H^q_{\mathfrak m}(M),E_R)
   \cong H^{-q}\mathbf R\mathrm{Hom}_{\widehat R}
                             (\widehat M,D_{\widehat R})
\]
of \(\widehat R\)-modules. If \(R\) is complete Cohen--Macaulay
of dimension \(e\), this becomes
\(\mathrm{Ext}^{e-q}_R(M,\omega_R)\), where
\(\omega_R=H^{-e}(D_R)\).
\end{corollary}
\begin{proof}
The local cohomology groups are \(\mathfrak m\)-power torsion
and are unchanged on completing \(M\), by \cite[L02, L09]{ca-topics}.
Every \(R\)-linear map between \(\mathfrak m\)-power torsion
modules is \(\widehat R\)-linear. This follows by using
\(R/\mathfrak m^r=\widehat R/(\mathfrak m\widehat R)^r\)
on submodules annihilated by \(\mathfrak m^r\). The
\(\widehat R\)-injective envelope of \(k\) restricts
to the \(R\)-injective envelope, as in \cite[M06]{ca-topics}.
Thus the first isomorphism follows by applying
Theorem~\ref{res:U19} to \(\widehat R\). For the second,
Theorem~\ref{res:U19} with \(M=R\) and the
Cohen--Macaulay cohomology criterion \cite[V07]{ca-topics}
show that \(D_R\) has a single nonzero cohomology
module in degree \(-e\). The module is the canonical
module by \cite[D08]{ca-topics}, and good truncation
identifies \(D_R\) with \(\omega_R[e]\). Now the
cohomology shift gives the claimed Ext group. These
are the completions and shifts of \cite[D07--D09]{ca-topics}.
\end{proof}

\section{An explicit projective case}

\begin{theorem}[Laurent monomials compute the twists on projective space]\label{res:U21}
Let \(k\) be a field, \(n\geq1\), \(X=\mathbf P^n_k\),
and \(S=k[x_0,\ldots,x_n]\) with its usual grading.
For \(m\in\mathbb Z\), the cohomology of \(\mathcal O_X(m)\)
is \(S_m\) in degree zero, vanishes for \(0<q<n\)
and for \(q>n\), and has basis
\[
   \{x_0^{a_0}\cdots x_n^{a_n}:\ a_i<0,\
           a_0+\cdots+a_n=m\}
\]
in degree \(n\).
\end{theorem}
\begin{proof}
Use the cover \(D_+(x_i)\). Its nonempty intersections
are affine Noetherian schemes, since the one-variable
dehomogenizations on a chart are polynomial rings and
localizations. Sections of \(\mathcal O(m)\) on the
intersection indexed by a nonempty subset \(J\) are the
degree-\(m\) part of \(S[x_j^{-1}:j\in J]\), by the
construction of the twisting sheaf. Quasi-coherent
modules on these affine intersections have vanishing
positive cohomology by the affine acyclicity argument of
\cite[L10]{ca-topics}. The \v{C}ech comparison
\cite[H02--H04]{ha-core} thus computes the cohomology.

In the \v{C}ech complex, decompose into exponent vectors
\(\mathbf a=(a_0,\ldots,a_n)\) of total degree \(m\).
The term for a subset \(J\) has a basis vector for
\(\mathbf a\) precisely when it contains the set
\(N=\{i:a_i<0\}\). Restrictions either carry the basis
vector to itself or have zero source, so each vector
gives a subcomplex. When \(N\) is nonempty and proper,
choose \(j\notin N\). The alternating simplex homotopy
pairs the components \(J\) not containing \(j\) with
\(J\cup\{j\}\). Explicitly the degree-minus-one map
takes the component at \(J\cup\{j\}\) to the component
at \(J\) with sign \((-1)^{|\{i\in J:i<j\}|}\), and has
zero component at \(J\) if \(j\in J\).
Both paired sets are admissible, because \(J\) contains
\(N\) and is therefore nonempty. Expanding the
alternating differential gives the identity on the
insertion of \(j\), while insertions of indices other
than \(j\) occur twice with opposite signs. Thus this
summand is acyclic. For \(N=\varnothing\), the same
homotopy works after adding the empty subset in
degree \(-1\), leaving exactly one copy of \(k\) in
degree zero of the unaugmented subcomplex.
For \(N\) the full set only the top-degree term remains.
Degreewise direct sums over exponent vectors preserve
cohomology. This proves the formulas.
\end{proof}

\begin{corollary}[Trace and duality for the twists]\label{res:U22}
In the situation of Theorem~\ref{res:U21}, the coefficient
of \((x_0\cdots x_n)^{-1}\) specifies a trace
\(H^n(X,\mathcal O_X(-n-1))\to k\).
Multiplication followed by trace gives perfect pairings
\[
 H^q(X,\mathcal O_X(m))\ \otimes_k\
 H^{n-q}(X,\mathcal O_X(-m-n-1))\ \longrightarrow k
 \qquad (0\leq q\leq n).
\]
The sheaf \(\mathcal O_X(-n-1)\) identifies with the
top differential forms \(\Omega^n_{X/k}\).
\end{corollary}
\begin{proof}
The trace space is one-dimensional by
Theorem~\ref{res:U21}. The only potentially nonzero
pairings are \(q=0,n\). In each of them multiplication
by a global section acts on the Laurent \v{C}ech
representative by ordinary multiplication. If \(x^{\mathbf a}\)
has all exponents nonnegative, its unique complementary
basis vector has exponents \(-1-a_i\) at each coordinate.
Their product gives coefficient \(1\), while any other
basis product gives coefficient zero. Thus the bases in
the two factors are dual. The formula also covers degrees
for which both spaces are zero.

For the top forms, on \(U_i=D_+(x_i)\) write
\(t_a=x_a/x_i\) and use the generator
\[
 \eta_i=(-1)^i\,d t_0\wedge\cdots
                   \wedge\widehat{d t_i}\wedge\cdots\wedge d t_n .
\]
For \(j\ne i\) put \(z=t_j=x_j/x_i\), a unit on the overlap.
The new coordinate indexed by \(i\) is \(1/z\), with
differential \(-z^{-2}dz\). Every other coordinate indexed
by \(a\ne i,j\) is \(t_a/z\), with differential
\(z^{-1}d t_a-z^{-2}t_a\,dz\). In wedging the new
differentials the summands with a second \(dz\) vanish.
The scalar factor without its sign is \(z^{-n-1}\), and
the sign of \(-d z\) moved from the \(i\)-slot to the \(j\)-slot
is \((-1)^{i+j}\). The factors \((-1)^i,(-1)^j\)
in the two generators cancel that sign. Hence
\(\eta_j=z^{-n-1}\eta_i\), the transition function of
the frame \(x_i^{-n-1}\) for \(\mathcal O(-n-1)\).
This gives the asserted identification.
\end{proof}
